\documentclass[aspectratio=169, 12pt]{beamer}

%--- Packages for Arabic and Layout ---%
\usepackage{polyglossia}
\setmainlanguage{arabic}
\setotherlanguage{english}

%--- Font Setup ---%
\newfontfamily\arabicfont[Script=Arabic, Scale=1.2]{Amiri}
\newfontfamily\quranfont[Script=Arabic, Scale=1.4]{Amiri}

%--- Theme & Color Setup ---%
\usetheme{Berlin}        
\usecolortheme{seahorse} 
\usefonttheme{serif}

%--- VISUAL TWEAKS: Vertical Space Optimization ---%
\setbeamertemplate{headline}{%
    \begin{beamercolorbox}[ht=2.5ex,dp=1.125ex]{section in head/foot}
        \insertsectionnavigationhorizontal{\paperwidth}{}{\hskip0pt plus1filll}
    \end{beamercolorbox}%
}

\setbeamertemplate{footline}{%
    \begin{beamercolorbox}[ht=2.5ex,dp=1.125ex,leftskip=.3cm,rightskip=.3cm plus1fil]{author in head/foot}
        \usebeamerfont{author in head/foot}\insertshortauthor \hfill \insertshortinstitute
    \end{beamercolorbox}%
}

\setbeamertemplate{section in head/foot}{\insertsectionhead\hspace{20pt}}


\setbeamertemplate{navigation symbols}{} 

%--- Metadata ---%
\title[أحكام التجويد]{ملخص أحكام تجويد القرآن الكريم}
\subtitle{شرح شامل للمخارج، الصفات، والمدود}
\author{دورة التجويد المكثفة المختصرة}
\institute{مبادئ في علم التجويد}
\date{\today}

\begin{document}

%--- Title Slide ---%
\begin{frame}
    \titlepage
\end{frame}

%--- Table of Contents (Right Aligned) ---%
\begin{frame}{فهرس المحتويات}
    \begin{columns}[t]
        % العمود الأيمن (يظهر أولاً في العربية)
        \column{0.45\textwidth}
        \raggedright % <--- هذا الأمر يلغي التمدد ويضبط المحاذاة لليمين
        \tableofcontents[sections={1-3}]
        
        % العمود الأيسر (يظهر ثانياً في العربية)
        \column{0.45\textwidth}
        \raggedright % <--- هذا الأمر يلغي التمدد ويضبط المحاذاة لليمين
        \tableofcontents[sections={4-7}]
    \end{columns}
\end{frame}

%===========================================================
% Section 1: Introduction
%===========================================================
\section{المقدمة}
\subsection{التعريف والحكم}

\begin{frame}{تعريف التجويد وحكمه}
    \begin{block}{تعريف التجويد}
        \begin{itemize}
            \item \textbf{لغة:} التحسين.
            \item \textbf{اصطلاحاً:} إخراج كل حرف من مخرجه مع إعطائه حقه (الصفات الذاتية) ومستحقه (الصفات العرضية).
        \end{itemize}
    \end{block}

    \begin{block}{حكم العمل به}
        \begin{itemize}
            \item \textbf{فرض عين:} فيما يغير المبنى أو يفسد المعنى (اللحن الجلي).
            \item \textbf{فرض كفاية:} فيما يتعلق بمهارة تحقيق الصفات الدقيقة.
        \end{itemize}
    \end{block}
\end{frame}

\subsection{المراتب والاستعاذة}
\begin{frame}{مراتب القراءة والاستعاذة}
    \begin{columns}
        \column{0.5\textwidth}
        \textbf{مراتب القراءة:}
        \begin{enumerate}
            \item \textbf{التحقيق:} القراءة باطمئنان وتؤدة (للتعليم).
            \item \textbf{التدوير:} مرتبة متوسطة بين التحقيق والحدر.
            \item \textbf{الحدر:} القراءة بسرعة مع مراعاة الأحكام.
        \end{enumerate}

        \column{0.5\textwidth}
        \textbf{أوجه الاستعاذة والبسملة:}
        \begin{itemize}
            \item قطع الجميع.
            \item وصل الجميع.
            \item قطع الاستعاذة ووصل البسملة بأول السورة.
            \item وصل الاستعاذة بالبسملة وقطعها عن أول السورة.
        \end{itemize}
    \end{columns}
\end{frame}

%===========================================================
% Section 2: Noon Sakinah
%===========================================================
\section{النون الساكنة}
\subsection{التعريف}

\begin{frame}{تعريف النون الساكنة والتنوين}
    \begin{itemize}
        \item \textbf{النون الساكنة:} نون لا حركة لها، تثبت لفظاً وخطاً.
        \item \textbf{التنوين:} نون ساكنة زائدة تلحق آخر الأسماء لفظاً لا خطاً.
    \end{itemize}
    \begin{alertblock}{الأحكام الأربعة}
        للنون الساكنة والتنوين أربعة أحكام: الإظهار، الإدغام، الإقلاب، الإخفاء.
    \end{alertblock}
\end{frame}

\subsection{الإظهار والإقلاب}
\begin{frame}{1. الإظهار الحلقي و 2. الإقلاب}
    \begin{columns}
        \column{0.5\textwidth}
        \textbf{الإظهار الحلقي}
        \begin{itemize}
            \item \textbf{حروفه:} (ء، هـ، ع، ح، غ، خ).
            \item مجموعة في: "أخي هاك علماً حازه غير خاسر".
            \item \textbf{الحكم:} إخراج النون واضحة بغير غنة.
        \end{itemize}

        \column{0.5\textwidth}
        \textbf{الإقلاب}
        \begin{itemize}
            \item \textbf{حروفه:} حرف الباء (ب).
            \item \textbf{الحكم:} قلب النون ميماً مخفاة مع الغنة.
            \item \textbf{مثال:} \quranfont{مِنْ بَعْدِ} تقرأ (ميم بعد).
        \end{itemize}
    \end{columns}
\end{frame}

\subsection{الإدغام والإخفاء}
\begin{frame}{3. الإدغام}
    \textbf{تعريفه:} إدخال حرف ساكن في متحرك بحيث يصيران حرفاً واحداً مشدداً.
    \vspace{0.3cm}
    \begin{table}
        \centering
        \begin{tabular}{|c|c|c|}
            \hline
            \textbf{النوع} & \textbf{الحروف} & \textbf{المجموع} \\
            \hline
            إدغام بغنة & (ي، ن، م، و) & ينمو \\
            \hline
            إدغام بغير غنة & (ل، ر) & -- \\
            \hline
        \end{tabular}
    \end{table}
    \textbf{مجموع الحروف:} (يرملون).
    \\ \small{* ملاحظة: الإدغام لا يكون إلا في كلمتين. إذا اجتمع في كلمة واحدة (مثل: صنوان) وجب الإظهار.}
\end{frame}

\begin{frame}{4. الإخفاء الحقيقي}
    \begin{itemize}
        \item \textbf{تعريفه:} النطق بالحرف بصفة بين الإظهار والإدغام مع بقاء الغنة.
        \item \textbf{حروفه:} (15 حرفاً) وهي الباقية.
        \item \textbf{مجموعة في أوائل البيت:}
        \\ \quranfont{صِفْ ذَا ثَنَا كَمْ جَادَ شَخْصٌ قَدْ سَمَا ... دُمْ طَيِّباً زِدْ فِي تُقًى ضَعْ ظَالِماً}.
    \end{itemize}
\end{frame}

%===========================================================
% Section 3: Mim Sakinah (Split into 2 Columns)
%===========================================================
\section{الميم الساكنة}
\subsection{أحكام الميم}

\begin{frame}{أحكام الميم الساكنة}
    للميم الساكنة ثلاثة أحكام:
    \vspace{0.3cm}
    \begin{columns}[t]
        % Column 1: Rules 1 & 2
        \column{0.5\textwidth}
        \textbf{1. الإخفاء الشفوي:}
        \begin{itemize}
            \item حرفه: الباء (ب).
            \item مثال: \quranfont{تَرْمِيهِم بِحِجَارَةٍ}.
        \end{itemize}
        \vspace{0.3cm}
        \textbf{2. الإدغام الشفوي:}
        \begin{itemize}
            \item حرفه: الميم (م).
            \item مثال: \quranfont{لَهُم مَّا يَشَاءُونَ}.
        \end{itemize}

        % Column 2: Rule 3
        \column{0.5\textwidth}
        \textbf{3. الإظهار الشفوي:}
        \begin{itemize}
            \item حروفه: باقي الحروف (26 حرفاً).
            \item \textbf{تنبيه:} أحذر من الإخفاء عند (الواو) و(الفاء)؛ لقرب المخرج.
        \end{itemize}
    \end{columns}
\end{frame}

\subsection{المشددتان}
\begin{frame}{أحكام النون والميم المشددتين}
    \begin{block}{القاعدة}
        يجب إظهار الغنة في النون والميم المشددتين بمقدار حركتين.
    \end{block}
    \textbf{أمثلة:}
    \begin{itemize}
        \item النون المشددة: \quranfont{إِنَّ}.
        \item الميم المشددة: \quranfont{ثُمَّ}.
    \end{itemize}
\end{frame}

%===========================================================
% Section 4: Madd
%===========================================================
\section{المدود}
\subsection{الأصلي والفرعي}

\begin{frame}{تعريف المد وأقسامه}
    \textbf{حروف المد:} الألف الساكنة المفتوح ما قبلها، الواو الساكنة المضموم ما قبلها، الياء الساكنة المكسور ما قبلها.

    \begin{columns}
        \column{0.5\textwidth}
        \textbf{1. المد الأصلي (الطبيعي):}
        ما لا تقوم ذات الحرف إلا به. مقداره حركتان.

        \column{0.5\textwidth}
        \textbf{2. المد الفرعي:}
        ما زاد عن الطبيعي لسبب (همز أو سكون).
    \end{columns}
\end{frame}

\begin{frame}{أنواع المد الفرعي (بسبب الهمز)}
    \begin{enumerate}
        \item \textbf{المد الواجب المتصل:} الهمزة بعد حرف المد في كلمة واحدة. (4-5 حركات).
        \\ مثال: \quranfont{السَّمَاء}.
        \item \textbf{المد الجائز المنفصل:} حرف المد في آخر كلمة والهمزة في أول التالية. (2، 4، 5 حركات).
        \\ مثال: \quranfont{إِنَّا أَعْطَيْنَاكَ}.
    \end{enumerate}
\end{frame}

\subsection{بسبب السكون}
\begin{frame}{أنواع المد الفرعي (بسبب السكون)}
    \begin{enumerate}
        \item \textbf{المد العارض للسكون:} أن يأتي بعد حرف المد سكون عارض للوقف. (2، 4، 6 حركات).
        \\ مثال: \quranfont{الْعَالَمِينَ} (عند الوقف).
        \item \textbf{المد اللازم:} أن يأتي بعد حرف المد سكون أصلي وصلاً ووقفاً. (6 حركات لزوماً).
        \begin{itemize}
            \item \textbf{كلمي:} (مثقل: \quranfont{الضَّالِّينَ} / مخفف: \quranfont{آلْآنَ}).
            \item \textbf{حرفي:} في فواتح السور (مجموعة في: كم عسل نقص).
        \end{itemize}
    \end{enumerate}
\end{frame}

%===========================================================
% Section 5: Makharij & Sifat
%===========================================================
\section{المخارج والصفات}
\subsection{المخارج}

% --- شريحة مخارج الحروف مع الصورة ---
\begin{frame}{المخارج العامة للحروف}
    \begin{columns}
        % العمود الأيمن للنص
        \column{0.45\textwidth}
        \textbf{المخارج العامة الخمسة:}
        \begin{enumerate}
            \item \textbf{الجوف:} خلاء الفم والحلق.
            \item \textbf{الحلق:} (أقصى، وسط، أدنى).
            \item \textbf{اللسان:} فيه 10 مخارج.
            \item \textbf{الشفتان:} (ف، م، ب، و).
            \item \textbf{الخيشوم:} للغنة.
        \end{enumerate}

        % العمود الأيسر للصورة
        \column{0.55\textwidth}
        \centering
        % تأكد من أن اسم الصورة مطابق تماماً (makharej.jpg) وموجودة في نفس المجلد
        \includegraphics[width=1\linewidth, height=0.8\textheight, keepaspectratio]{makharej.jpg}
    \end{columns}
\end{frame}

\subsection{الصفات}
\begin{frame}{صفات الحروف (التي لها ضد)}
    \begin{table}
        \centering
        \scriptsize
        \begin{tabular}{|c|c|c|}
            \hline
            \textbf{الصفة} & \textbf{ضدها} & \textbf{حروف الصفة الأولى} \\
            \hline
            الهمس & الجهر & فحثه شخص سكت \\
            \hline
            الشدة & الرخاوة & أجد قط بكت \\
            \hline
            الاستعلاء & الاستفال & خص ضغط قظ \\
            \hline
            الإطباق & الانفتاح & ص، ض، ط، ظ \\
            \hline
        \end{tabular}
    \end{table}
\end{frame}

\begin{frame}{صفات الحروف (التي لا ضد لها)}
    \begin{itemize}
        \item \textbf{القلقلة:} اضطراب المخرج (قطب جد).
        \item \textbf{الصفير:} (ص، ز، س).
        \item \textbf{اللين:} (و، ي) الساكنتين المفتوح ما قبلهما.
        \item \textbf{الانحراف:} (ل، ر).
        \item \textbf{التكرير:} (ر).
        \item \textbf{التفشي:} (ش).
        \item \textbf{الاستطالة:} (ض).
    \end{itemize}
\end{frame}

%===========================================================
% Section 6: Miscellaneous
%===========================================================
\section{أحكام متفرقة}
\subsection{التفخيم والراء}

\begin{frame}{التفخيم والترقيق}
    \begin{enumerate}
        \item \textbf{مفخم دائماً:} حروف الاستعلاء (خص ضغط قظ).
        \item \textbf{مرقق دائماً:} باقي الحروف.
        \item \textbf{دائر بينهما:} الألف، لام لفظ الجلالة، الراء.
    \end{enumerate}
    \begin{block}{أحكام الراء}
        \textbf{تفخم:} إذا كانت مفتوحة أو مضمومة. \\
        \textbf{ترقق:} إذا كانت مكسورة.
    \end{block}
\end{frame}

\subsection{الوقف}
\begin{frame}{أقسام الوقف}
    \begin{enumerate}
        \item \textbf{التام:} تم معناه ولم يتعلق بما بعده.
        \item \textbf{الكافي:} تعلق بما بعده معنى لا لفظاً.
        \item \textbf{الحسن:} تعلق بما بعده لفظاً ومعنى.
        \item \textbf{القبيح:} لم يتم معناه أو أوهم معنى فاسداً.
    \end{enumerate}
\end{frame}

%--- End Slide ---%
\begin{frame}
    \centering
    \Huge سبحانك اللهم وبحمدك، نشهد أن لا إله إلا أنت، نستغفرك ونتوب إليك.
\end{frame}

\end{document}
